% !PRINTED 686
\weberpage{686}
%<<CONT
is a root of unity. The system of exponents of the reduced number
$\alpha_0$ derived from $\alpha$ is completely determined by $\alpha$ itself.
This proves the theorem:

II. \gesperrt{To every number $\alpha$ (different from zero) of the
field $\Omega$ there are $w$, and not more, reduced
numbers $\varrho\,\alpha_0$.}

If in place of the fundamental system $(\varepsilon_1, \varepsilon_2, \ldots, \varepsilon_{\nu-1})$ another
one is used, the exponents $\xi_1, \xi_2, \ldots, \xi_{\nu-1}$ undergo
a linear substitution with integral coefficients of determinant $\pm 1$.
A reduced number may then cease to be reduced for the new fundamental
system. The concept of the reduced number
therefore depends on the choice of the fundamental system. But the
number of the reduced numbers derived from a given number
is always the same.

By this theorem we have attained our purpose, namely to single
out, from the whole infinite system of numbers associated with one
another, a definite finite number of representatives.

\begin{center}
\S.\,189.\\[4pt]
\gesperrt{Limits of the number of the integers of the field $\Omega$
divisible by an ideal.}
\end{center}
% [source note] The following finiteness assertion and definition of T omit counting modulo association. Infinitely many units already have absolute norm 1. The intended finite count is of non-associated integers (equivalently the corresponding principal ideals), explicitly specified on printed p. 690. The source also changes the endpoint from norm at most t here to natural numbers less than t below; this can alter a boundary term but not the limiting quotient. The wording is preserved.

Our earlier considerations have shown that every
integer of a field $\Omega$ has only a finite number of ideals
as divisors. Since every integral ideal is a divisor of its norm,
there are therefore also only a finite number of integers
in $\Omega$ whose absolute norm does not exceed a given
limit.

From these numbers we now single out a part once more,
and ask for the number $T$ of all the \gesperrt{integers}
of the field $\Omega$ whose absolute norm does not exceed a positive quantity $t$
and which are divisible by a given ideal $\mathfrak{a}$.

Let us take the rational field as a preliminary example;
there $T$ is the number of the natural numbers that are less than
$t$ and divisible by a given integer $m$, that is, in
the notation used earlier (\S.\,182), $T = E\left(\frac{t}{m}\right)$.
